\documentclass[12pt,letterpaper]{report}
\usepackage{graphicx}
\usepackage{scrextend}
\usepackage{vmargin}
\usepackage{graphicx}
\usepackage{multirow}
\usepackage[utf8]{inputenc}
\usepackage[spanish]{babel}
\usepackage{multicol}
\usepackage{enumerate}
\usepackage{float}
\usepackage{amsmath, amsthm, amssymb, amsfonts}
\usepackage[usenames]{color}
\parindent=0mm
\pagestyle{empty}
\definecolor{miorange}{rgb}{0.91, 0.43, 0.0}
\begin{document}
\setmargins{2.5cm}      
{1.5cm}                     
{2cm}  
{24cm}                    
{10pt}                          
{1cm}                          
{0pt}                             
{2cm}
\begin{flushright}
\today
\end{flushright}
\vspace{-1.75cm}
\section*{Tarea Clase 8}
\subsection*{Obtenga las soluciones de las siguientes ecuaciones:}
\begin{enumerate}
    \item Hemos comprado 18L de pintura en una tienda de bricolaje donde el precio de la pintura azul es 
    12\$/Ly el de la pintura verde es 13.5\$/L. ¿Cuántos litros de pintura de cada color hemos comprado gastando 
    234\$?
    \item La abuela de Pedro quiere dar dinero a sus nietos para las vacaciones de Navidad. Si le da 
    25\$ a cada uno, le sobrarían 25\$. Sin embargo, si les diera 35
    \$, le faltarían 25\$. ¿De cuánto dinero dispone la abuela de Pedro?
    \item Hace 5 años, la edad de Sonia era triple que la de Roberto, y dentro de 10 años será
    doble. ¿Qué edad tiene cada uno? 
    \item En un test de elección múltiple, se puntúa 4 por cada respuesta correcta y se resta un
    punto por una equivocada. Un estudiante responde a 17 cuestiones y obtiene 43
    puntos. ¿Cuántas cuestiones respondió correctamente?
    \item En un aula, la asignatura de gimnasia la han aprobado el 62,5\% de las alumnas y el 80\% de los alumnos, mientras que la asignatura de historia la han aprobado 87,5\% de las alumnas y el 60\ de los alumnos: 
    \begin{table}[H]
        \centering 
        \begin{tabular}{|c|c|c|}\hline
            & \textbf{Gimnasia} & \textbf{Historia} \\ \hline
            Alumnas & 62.5\% & 87.5\% \\ \hline
            Alumnos & 80\% & 60\% \\ \hline
            \textbf{Total} & \textbf{26} & \textbf{26}\\ \hline
        \end{tabular}
    \end{table}
    \item Un avión dispone de 32 asientos en clase A y de 50 asientos en clase B cuya venta supone un total de 14.600€. Sin embargo, sólo sea han vendido 10 asientos en clase A y 40 en clase B, obteniendo un total de 7.000€.¿Cuál es precio de un asiento en cada clase?
    \item Hallar un número de dos cifras sabiendo que la suma de las cifras es 12 y que la primera de ellas es el triple de la segunda.
\end{enumerate}
Las respuestas de las tareas pueden ser enviadas al siguiente correo giovannilopez9808@gmail.com con el asunto entre corchetes de la siguiente manera [Tarea \#] y el nombre del archivo con sus apellidos y nombre, ya sea el nombre completo o solo una parte, por ejemplo \textit{LopezPadilla.pdf}, \textit{LopezGiovanni.pdf}, \textit{LopezPadillaGiovanniGamaliel.pdf}. Las respuestas envienlas en formato pdf, estas pueden ser totalmente digitales (que escriban todo en la computadora) o que sean fotografías, pero por favor, envielo en solo un documento.
\end{document}